\begin{example}
    Las sucesiones $(an^\sigma\log^\tau n)_{n\geq2}$ con $a \neq 0$, 
    $\sigma \in (0,1)$ y $\tau \in \mathbb{R}$ son equidistribuidas módulo $1$.
\end{example}

\begin{resolve}
    Sea $x_n = an^\sigma\log^\tau n$ para cada $n \geq 2$. Definimos $g(x) = a(x+1)^\sigma\log^\tau (x+1)$
    en~$[1,\infty)$, entonces $g(n) = x_{n+1}$ para cada $n \geq 1$, y su derivada es 
    \[
        g'(x) = a (x+1)^{\sigma-1} \log^{\tau-1}(x+1) \left(\sigma \log(x+1) + \tau\right)
    \]
    esta es continua en $[1,\infty)$ por ser composición de continuas y como $\sigma \in (0,1)$, 
    \[
        \lim_{x \to \infty} g'(x) = 0
        \qquad \text{y} \qquad
        \lim_{x \to \infty} x |g'(x)| = \infty
    \]

    Monotonia, estudiar como hacerlo.

    Por el teorema \ref{teo:fejer_euler}, $(g(n))$ es equidistribuida módulo $1$ y 
    por lo tanto, $(x_n)$ también lo es.
\end{resolve}